% Generated draft: requires mathematical and source review.
\documentclass[11pt]{article}
\usepackage[margin=25mm]{geometry}
\usepackage{fontspec}
\setmainfont{Latin Modern Roman}
\usepackage{amsmath,amssymb,mathtools,unicode-math}
\setmathfont{Latin Modern Math}
\usepackage{enumitem,needspace,xurl,hyperref}
\setlength{\parindent}{0pt}
\setlength{\parskip}{5pt}
\setlength{\emergencystretch}{3em}
\newcommand{\N}{\mathbb{N}}
\newcommand{\Z}{\mathbb{Z}}
\newcommand{\Q}{\mathbb{Q}}
\newcommand{\R}{\mathbb{R}}
\newcommand{\C}{\mathbb{C}}
\newcommand{\F}{\mathbb{F}}
\newcommand{\A}{\mathbb{A}}
\newcommand{\PP}{\mathbb{P}}
\newcommand{\E}{\mathbb{E}}
\newcommand{\eps}{\varepsilon}
\newcommand{\dd}{\,\mathrm{d}}
\begin{document}
\section*{UnsolvedMath ID: 30000028}
\textbf{Source problem code:} catalogue entry
\subsection{Definitions and mathematical statement}
The source asks whether equivalence classes of global deformations of a Lie algebra carry a natural topological or geometric moduli space. The supplied literature assessment specifically describes a broad infinite-dimensional question and progress for Virasoro-type algebras; the main extracted sentence does not specify the Lie algebra, coefficient field, base category, or equivalence relation.

A formulation must first fix these data. For instance, after choosing a field $k$ and a Lie algebra $\mathfrak g$, a global deformation over a pointed base $(B,b)$ can mean a flat family $\mathcal L$ of $\mathcal O_B$-Lie algebras, together with a marking
\[
 \mathcal L\otimes_{\mathcal O_B}k(b)\simeq\mathfrak g.
\]
Over the same base, marked equivalence means a bracket-preserving $\mathcal O_B$-linear isomorphism compatible with the marking. Pullback along pointed base maps then gives a deformation functor, or a groupoid-valued deformation problem retaining automorphisms. In finite rank one may impose local freeness; for an infinite-dimensional algebra the appropriate module class, grading, topology and continuity conditions must instead be specified. These are alternatives to be fixed, not extra hypotheses silently supplied by the record.

The research question is to construct a moduli object naturally encoding the chosen global families and their equivalences. A fine moduli space would require a universal-family property; a coarse space, stack, or topological quotient has a different universal property. Merely placing an arbitrary topology on a set of classes does not resolve the intended ``natural'' construction. Global families over actual bases must also be distinguished from formal deformations over completed local rings or infinitesimal Artinian bases. The source does not say whether classes of families over different bases, classes of fibers, or one fixed-base deformation functor are to be classified, so no single representability theorem is asserted here.

The supplied literature triage checked 21 August 2026 reports nontrivial global families and geometric moduli methods for Virasoro-type algebras, but its reviewer did not locate a complete classification for the broad infinite-dimensional question. That bounded literature assessment is not a proof of nonexistence of a moduli space or an independent current survey.
\subsection{Short English statement}
Can global Lie-algebra deformations be organized into a natural moduli space? The question needs a class of Lie algebras, allowed bases and families, equivalences and a moduli universal property. The supplied assessment reports partial progress for Virasoro-type algebras, without a complete answer to the broad infinite-dimensional formulation.
\subsection{Sources}
\begin{itemize}
\item Local UnsolvedMath registry, permanent ID 30000028, exact source code catalogue entry; \emph{Moduli Spaces of Global Lie-Algebra Deformations}. Proposed year 2004, status partially\_solved. The supplied corrected\_verified wording asks about a natural topological or geometric moduli space.
\item \emph{Finite and Infinite Dimensional Complex Geometry and Representation Theory (2004)}, Oberwolfach Reports, DOI 10.4171/owr/2004/06, pp. 295--348, \url{https://doi.org/10.4171/owr/2004/06}.
\item Supplied literature triage checked 2026-08-21: ``Oberwolfach report'', \url{https://publications.mfo.de/bitstream/handle/mfo/2825/OWR_2004_06.pdf}; ``Fialowski--Schlichenmaier'', \url{https://arxiv.org/abs/math/0610851}; ``Deformation survey'', \url{https://www.renyi.hu/~fialowsk/pubs-af/BulletinTICMI2016.pdf}. This dated OWR-derived curation requests independent verification of time-sensitive status and cited claims; linked PDFs were not accessed here.
\end{itemize}
\end{document}
